\subsubsection{UseCases}
Bevor detailierte technische Anforderungen an das System ausgearbeitet werden koennen, benoetigt es eine einfache Uebersicht der Kernanwendungsfaelle.
Hierzu wurde eine UseCase-Diagramm erstellt, um diese stark abstrahiert und einfach darzustellen, sodass die Stakeholder des Projektes diese einfach verstehen und nachvollziehen koennen.
Das folgende Diagramm zeigt die zentralen Use-Cases bei der Nutzung eines Bewerbermanagementportals.
Die zwei wichtigsten Anwender sind in diesem Fall die Recruter-innen und die Bewerber-innen.
\begin{figure}[H]
    \centering
    \includegraphics[width=0.7\textwidth]{UseCase.jpg}
    \caption{Use-Cases}
    \label{fig:useCase}
\end{figure}
Fuer den Bewerber sind die notwendigsten Funktionalitaeten:
\begin{itemize}
\item alle ausgeschriebenen Stellen durchsuchen
\item einzelne Stellenanzeigen betrachten
\item Bewerbungen durchfuehren und
\item aktuellen Bewerbungsstatus einsehen
zu koennen.
\end{itemize}
Die bereitzustellenden Funktionalitaeten fuer den Recruiter sind im Gegensatz:
\begin{itemize}
\item Stellenanzeigen erstellen
\item Stellenanzeigen bearbeiten
\item Stellenanzeigen aktivieren/deaktivieren
\end{itemize}
zu koennen.
\newline
Einen kleinen spezifischen, aber sehr wichtigen, Anwendungsfall stellt das Vergeben der Anmeldedaten fuer den Recruiter dar.
Dies ist notwendig, da sich die Zugriffe zwischen Recruiter-innen und Bewerber-innen unterscheiden. 
